Form Wizard দিয়ে Student-এর প্রয়োজনীয় ফিল্ড নিয়ে একটি ফর্ম তৈরি করো। নতুন রেকর্ড এন্ট্রি করে টেবিলে তা এসেছে কি না দেখো। শিক্ষার্থীর সঙ্গে ফল দেখাতে মূল ফর্ম ও Result সাবফর্ম StudentID দিয়ে যুক্ত করো। ভিন্ন শিক্ষার্থী বাছলে উপফর্মে কেবল সংশ্লিষ্ট ফল আসছে কি না পরীক্ষা করো।

Report Wizard-এ ফলাফলের কুয়েরি বেছে Roll, Name ও ICTMark নাও। প্রয়োজনমতো পরীক্ষাভিত্তিক group এবং রোলভিত্তিক sort দাও। Preview-তে সব কলাম, পৃষ্ঠা নম্বর ও শিরোনাম ঠিক আছে কি না যাচাই করো। একই তথ্য টেবিলে, ফর্মে ও রিপোর্টে আলাদাভাবে দেখালেও কোনটি মূল সংরক্ষণস্থান তা খাতায় ব্যাখ্যা করো।

\textbf{সুরক্ষার অনুশীলন:} অনুশীলনের ডেটাবেজের আলাদা কপি রেখে সমর্থিত Access ফরম্যাটে Encrypt with 